\documentclass[letterpaper]{article} % DO NOT CHANGE THIS
\usepackage[submission]{aaai2027}  % DO NOT CHANGE THIS
% The serif, sans-serif, and monospaced fonts are loaded automatically by
% aaai2027.sty (newtxtext, helvet, courier). DO NOT add \usepackage{times},
% \usepackage{helvet}, \usepackage{courier}, or any other font package.
\usepackage[hyphens]{url}  % DO NOT CHANGE THIS
\usepackage{graphicx} % DO NOT CHANGE THIS
\urlstyle{rm} % DO NOT CHANGE THIS
\def\UrlFont{\rm}  % DO NOT CHANGE THIS
\usepackage{natbib}  % DO NOT CHANGE THIS AND DO NOT ADD ANY OPTIONS TO IT
\usepackage{caption} % DO NOT CHANGE THIS AND DO NOT ADD ANY OPTIONS TO IT
\frenchspacing  % DO NOT CHANGE THIS
%
% These are recommended to typeset algorithms but not required.
\usepackage{algorithm}
\usepackage{algorithmic}
%
% Math support (standard, permitted).
\usepackage{amsmath}
\usepackage{amssymb}
%
% Recommended for better-looking tables
\usepackage{booktabs}
%
% Visible red placeholder for pending numeric results -- grep "\PH{" to find them.
\usepackage{xcolor}
\newcommand{\PH}[1]{{\color{red}\textbf{[#1]}}}
%
% Keep the \pdfinfo as shown here.
\pdfinfo{
/TemplateVersion (2027.1)
}

\setcounter{secnumdepth}{0} %May be changed to 1 or 2 if section numbers are desired.

% Title -- mixed (Title) case.
\title{PACE: Curating Failure Modes for Query-Efficient\\ Interactive Imitation Learning}

\author{
    Anonymous Submission
}
\affiliations{
    %
}

\begin{document}

\maketitle

\begin{abstract}
Interactive imitation learning corrects the compounding covariate shift of
behavior cloning by aggregating fresh expert demonstrations on visited states.
The dominant DAgger family answers only \emph{when} to query the expert---via a
per-state uncertainty or safety gate---leaving implicit \emph{which} of the many
observed failures deserve scarce expert effort. \textbf{PACE}
(\textbf{P}erceive~$\rightarrow$~\textbf{A}ssess~$\rightarrow$~%
\textbf{C}hoose~$\rightarrow$~\textbf{E}xecute) is an interactive-imitation loop
for diffusion policies that curates each round's rollout failures as a batch.
\emph{Perceive} anchors a geometric (optionally visual) descriptor at the
peak-diffusion-loss step; \emph{Assess} clusters descriptors into failure modes
scored by dominance and cross-round novelty; \emph{Choose} rotates the target via
cross-round memory to a mode not yet covered; \emph{Execute} collects the
correction on-policy at the failure, or optionally from an LLM-prescribed
sub-task-entry start state. PACE's central axis is thus \emph{which} failure mode
to correct, not only \emph{when} to query; the prescribed \emph{where}-to-start is
an ablated option, and a unifying view places query-based baselines in one corner
of the design space. We study PACE on two \emph{implemented} tasks (a $5\times5$
grid and ManiSkill Push, in state and image observations, five seeds) against the
DAgger family and single-axis \emph{where} and \emph{which} controls; three
RoboSuite tasks (Lift, Wipe, Door) are \emph{planned} extensions with no result
rows. Under a paired per-seed protocol we test whether PACE reaches the target
success rate with \PH{pace-vs-diff-q-reduction} fewer expert queries than the
strongest baseline. All reported numbers are placeholders pending the five-seed
runs; today's contribution is the method, framework, and pre-registered protocol.
\end{abstract}

\section{Introduction}

Behavior cloning regresses an expert's actions on the states the expert
visits~\citep{pomerleau1988alvinn,pomerleau1991efficient}, and its
well-known failure mode is \emph{covariate shift}: the learner's own
actions determine the states it later encounters, so small prediction
errors push it off the demonstrated support, where errors compound over
the horizon~\citep{ross2011dagger}. Interactive imitation learning (IIL)
confronts this by aggregating expert labels on states the \emph{learner}
visits, bounding error linearly rather than quadratically in the
horizon~\citep{ross2011dagger,celemin2022iil}. The price is expert
effort: vanilla DAgger relabels every rolled-out state, which is
prohibitive when the expert is a human teleoperator or an expensive
planner~\citep{zhang2017safedagger,hoque2021thriftydagger}. \emph{Query
efficiency} --- reaching a target success rate with as few expert queries
as possible --- is therefore the axis along which IIL methods are judged,
and the one we study.

\paragraph{The DAgger-family gap.}
Every query-efficient descendant of DAgger decides \emph{when} to hand
control to the expert during a rollout; only the signal varies. SafeDAgger
uses a learned safety classifier~\citep{zhang2017safedagger}; DropoutDAgger
and EnsembleDAgger read predictive uncertainty off dropout and ensemble
variance~\citep{menda2017dropoutdagger,menda2019ensembledagger,
gal2016dropout,lakshminarayanan2017ensembles}; HG-DAgger and LazyDAgger watch
the learner--expert discrepancy~\citep{kelly2019hgdagger,hoque2021lazydagger};
ThriftyDAgger budgets by novelty and risk~\citep{hoque2021thriftydagger}; and
Diff-DAgger --- built for the multimodal diffusion policies we adopt --- reuses
the diffusion training loss as the query signal~\citep{lee2025diffdagger,
chi2023diffusionpolicy}. Different signals, one decision structure: a per-state
predicate $\mathrm{Query}(s_t)$ fires the first time a scalar score crosses a
threshold, and the expert takes over from that step~(Eq.~\ref{eq:iil-loop}); they
differ only in the score inside the indicator
(Eqs.~\ref{eq:safe}--\ref{eq:diffdagger}). This answers only \emph{when} to query
and leaves two axes unasked. Nothing decides \emph{which} of a round's distinct
failures deserve correction, so redundant failures from one mode are corrected
repeatedly while rarer modes starve; and with no memory across rounds, a
mode already fixed can be re-queried indefinitely. Budget is thus spent on the
\emph{easiest-to-detect} uncertainty rather than a diverse cover of the policy's
failure modes. A third, softer axis --- \emph{where} a corrective demonstration
begins --- the baselines fix to whatever state the rollout reached; PACE can
prescribe this start state, but treats it as an option, not the core method.

\paragraph{PACE.}
\textbf{PACE} (\textbf{P}erceive $\to$ \textbf{A}ssess $\to$
\textbf{C}hoose $\to$ \textbf{E}xecute) widens the DAgger-family query from
``\emph{when}'' to ``\emph{when} $+$ \emph{which},'' with an \emph{optional}
\emph{where}. Each round it rolls the current diffusion policy to collect a
\emph{set} of failures rather than one intervention point; the four stages are
defined in \S\ref{sec:method}. Perceive localizes each failed rollout at its
peak diffusion-loss frame --- the state the Diff-DAgger signal
flags~(Eq.~\ref{eq:diffdagger}) --- as a geometric or visual
descriptor~(Eq.~\ref{eq:phi}) that a vision--language
model~\citep{bai2025qwen3vl} describes. Assess clusters descriptors into failure
\emph{modes}~\citep{lloyd1982kmeans}, sets their count by
silhouette~(Eq.~\ref{eq:partition}), and marks the dominant one. Choose is where
coverage accumulates: a cross-round memory penalizes covered modes and
\emph{rotates} the target away from solved ones~(Eq.~\ref{eq:memory}), then a
farthest-point / $k$-center rule fixes the anchor failure and a context
set~(Eq.~\ref{eq:diversity}). Execute defaults to an on-policy expert
demonstration from the anchor's own step; with the option on, a language
model~\citep{yang2025qwen3} prescribes a corrective sub-task-entry start state,
projects it onto the task's workspace bounds, and the expert rolls from
there~(Eqs.~\ref{eq:prescribe}--\ref{eq:prescribe-demo}). The primary novelty is
the \emph{which} axis (Perceive $+$ Assess $+$ Choose): curating a round's
failures into modes and spending the single query on a memory-rotated,
previously-uncorrected mode. The prescribed \emph{where} start state is an option
we expose and ablate. This decomposition also gives a unifying view in which the
query-based baselines occupy one corner of the PACE design space; we develop that
view in \S\ref{sec:method} rather than here.

Two lines of prior work meet here. The DAgger family supplies correction on the
learner's own state distribution~\citep{ross2011dagger,laskey2017dart}, whose
diffusion loss we reuse as the failure signal~\citep{lee2025diffdagger}.
Foundation-model robotics supplies the machinery to \emph{perceive} and describe
failures~\citep{liu2023reflect,duan2025aha,shinn2023reflexion} and to
\emph{prescribe} structured targets --- here reset states and sub-task entry
points, not whole plans or reward functions~\citep{liang2023codeaspolicies,
ma2024eureka,yu2023language2rewards,florensa2017reversecurriculum,
eysenbach2018leavenotrace}. Assess and Choose are, in effect, batch active
learning over failures: cluster into modes and cover them with a diversity-aware
coreset under a budget~\citep{settles2009active,ash2020badge,sener2018coreset}.
Prior work touches each axis in isolation --- reverse-curriculum methods choose
\emph{where} to start~\citep{florensa2017reversecurriculum,eysenbach2018leavenotrace},
intervention-reweighting up-weights \emph{which} samples
matter~\citep{mandlekar2020iwr,liu2023sirius}, and failure-reasoning VLMs perceive
and cluster failures~\citep{liu2023reflect,duan2025aha}. PACE's novelty is their
\emph{integration}: perceived failure-mode clustering feeds a memory-rotated
target inside one diffusion-policy DAgger loop under a fixed one-demo-per-round
budget, with selection made \emph{before} data is collected rather than after. To
separate the integration from its parts we run two single-axis controls alongside
the DAgger family: a reset-curriculum \emph{where}-arm~\citep{florensa2017reversecurriculum}
without failure-mode curation, and an IWR-style \emph{which}-arm~\citep{mandlekar2020iwr}
that up-weights failure states after collection. We do not claim to be first to
consider \emph{which} or \emph{where}; we claim this failure-mode-driven
combination and evaluate what it buys over each axis alone.

\paragraph{Evaluation.}
Our scope is two implemented tasks, each in a state and an image modality
(Table~\ref{tab:tasks}): the toy $5{\times}5$ grid (multimodal navigation,
shortest-path expert) and Push (ManiSkill
PushT~\citep{mu2021maniskill,tao2024maniskill3,florence2021implicitbc}, a
diffusion-policy learner~\citep{chi2023diffusionpolicy,ho2020ddpm}). We compare
against the query-efficient DAgger family, a uniform-random control, and the two
single-axis controls. Three RoboSuite/UR5 tasks --- Lift, Wipe,
Door~\citep{zhu2020robosuite,mandlekar2021robomimic} --- are \emph{planned}
extensions carrying no result rows. All robot methods share one diffusion-policy
learner and differ only in the demonstration-acquisition rule, so the comparison
isolates query efficiency. We measure held-out SR, expert queries to reach
$\mathrm{SR}_{\mathrm{target}}{=}0.90$~(Eqs.~\ref{eq:sr}--\ref{eq:queries}), demo
efficiency (Eq.~\ref{eq:eff}), and coverage~(Eq.~\ref{eq:cov}), reporting PACE's
extra reset/foundation-model cost separately (\S Cost Accounting) and the sign of
each comparison from paired per-seed differences over five seeds.

\begin{table}[t]
\centering
\small
\caption{The evaluation suite. Every task is run in both a state and an image
modality; robot tasks share a single diffusion-policy learner across all
methods so that only the demonstration-acquisition rule varies. $\dim(\mathcal{A})$
is the policy action space (T2 policy control is $9$-D \texttt{rel\_joint\_pos};
its PPO expert emits $7$ arm-joint deltas, which we lift to $9$-D by placing them
in the $7$ arm-joint components and setting the two \texttt{panda\_stick} finger
components --- which have no actuator on this end-effector --- to the novice's own
(zero-delta) values, so $\|\cdot\|$ between novice and expert is taken in the
shared $9$-D space with the two inert dimensions contributing zero to every
discrepancy identically for all methods --- see
\S Baselines). \textbf{Diff-DAgger} is a robot-only baseline (requires the
diffusion learner). \textbf{T1--T2 are implemented and are the only tasks
evaluated in this paper.} T3--T5 (RoboSuite/UR5) are a \emph{planned} protocol
only: not implemented, not run, and carrying no result rows anywhere in the paper.}
\label{tab:tasks}
\begin{tabular}{llcccc}
\toprule
 & Sim & Embodiment & $\dim(\mathcal{A})$ & Expert & Status \\
\midrule
T1 Toy  & MazeNav          & grid agent    & $4$ (disc.) & $A^{\!\star}$/BFS & done \\
T2 Push & ManiSkill3       & Panda stick   & $9$         & PPO & done \\
T3 Lift & RoboSuite        & UR5e          & TBD & MP & planned \\
T4 Wipe & RoboSuite        & UR5           & TBD & MP & planned \\
T5 Door & RoboSuite        & UR5           & TBD & MP & planned \\
\bottomrule
\end{tabular}
\end{table}

\paragraph{Contributions.}
We make three contributions; because the empirical numbers are pending,
\textbf{(3)} ships as a \emph{pre-registered protocol and code} rather than
finished results.
\textbf{(1)} We reframe interactive imitation learning
\cite{ross2011dagger,celemin2022iil} as a per-round failure-curation problem
rather than a per-state query-gating problem, and cast the safe /
query-efficient DAgger family---SafeDAgger, DropoutDAgger, EnsembleDAgger,
ThriftyDAgger, and the diffusion-native Diff-DAgger
\cite{zhang2017safedagger,menda2017dropoutdagger,menda2019ensembledagger,hoque2021thriftydagger,lee2025diffdagger}---%
as the special case that answers only \emph{when} to query
(Eq.~\ref{eq:iil-loop}, Eq.~\ref{eq:diffdagger}).
\textbf{(2)} \textbf{PACE} itself, a Perceive--Assess--Choose--Execute loop
for diffusion-policy learners \cite{chi2023diffusionpolicy,ho2020ddpm} whose
central added axis is \emph{which} failure mode to correct: it perceives failures
by peak-loss localization and (for prescription) a vision--language model
\cite{bai2025qwen3vl,liu2023reflect,duan2025aha}, partitions them into modes by
clustering \cite{lloyd1982kmeans}, and chooses a memory-rotated target with a
farthest-point / k-center context set \cite{sener2018coreset,eldar1997fps}; an
\emph{optional} Execute stage additionally prescribes an LLM-generated corrective
sub-task-entry start state
\cite{yang2025qwen3,florensa2017reversecurriculum,eysenbach2018leavenotrace}
(Eqs.~\ref{eq:perceive}--\ref{eq:prescribe-demo}). Its novelty is the specific
integration of these axes in one loop, not any single axis alone; to separate the
integration from its parts we include a reset-curriculum \emph{where}-arm and an
IWR-style \emph{which}-arm as run baselines.
\textbf{(3)} We contribute a \emph{controlled, pre-registered empirical protocol}
(and its implementation) on the two implemented tasks
(toy grid, ManiSkill Push) in both state and image modalities
\cite{tao2024maniskill3,florence2021implicitbc,nair2022r3m}: five seeds, a paired
per-seed test, up to eight demonstration-acquisition arms (Diff-DAgger on robot
tasks only) under a matched one-demo-per-round budget, and a per-cell,
effect-size ablation that isolates each PACE stage and the foundation models.
The protocol is designed to test whether PACE cuts
expert queries to reach $\mathrm{SR}_{\mathrm{target}}{=}0.90$ by
\PH{pace-vs-diff-q-reduction} versus the best baseline at a mean held-out success
rate of \PH{mean-pace-sr}. \emph{Every result
cell is a placeholder pending the five-seed runs; no numeric claim in this paper
is yet established, and we phrase the empirical section as the protocol that will
answer it. What we can defend now is the method, the framework, and the protocol.}

\section{Related Work}
\label{sec:related}

Interactive-IL baselines answer a single question through a per-state predicate
$\mathrm{Query}(s_t)\!=\!\mathbf{1}[\,\mathrm{score}(s_t)\,\square\,\mathrm{thresh}\,]$
(Eqs.~\eqref{eq:iil-loop}, \eqref{eq:safe}--\eqref{eq:diffdagger}); PACE instead
perceives and partitions a round's failures, chooses a memory-rotated mode to
correct, and executes the correction (on-policy, optionally at a prescribed entry
state) (Eqs.~\eqref{eq:perceive}--\eqref{eq:prescribe-demo}). We review each thread
PACE touches, marking the axis it adds.

\subsection{Interactive Imitation Learning and the DAgger Family}
Behavior cloning maps observations to expert actions, dating to ALVINN
\cite{pomerleau1988alvinn,pomerleau1991efficient}, but suffers covariate shift:
small errors compound into states absent from $\mathcal{D}$
(Eqs.~\eqref{eq:expert}--\eqref{eq:bc}) \cite{ross2011dagger,laskey2017dart}. DAgger
recasts imitation as no-regret online learning, aggregating expert labels on
\emph{learner}-visited states to bound error linearly rather than quadratically
\cite{ross2011dagger}; AggreVaTe and Deeply AggreVaTeD extend this to cost-to-go and
differentiable policies \cite{ross2014aggrevate,sun2017aggrevated}. Since vanilla
DAgger demands
near-constant expert access, \emph{safe}, query-efficient variants gate \emph{when}
to defer: SafeDAgger learns a discrepancy classifier (Eq.~\eqref{eq:safe})
\cite{zhang2017safedagger}, and DropoutDAgger/EnsembleDAgger threshold
Bayesian-dropout and ensemble-variance signals
(Eqs.~\eqref{eq:dropout}--\eqref{eq:ensemble})
\cite{menda2017dropoutdagger,menda2019ensembledagger}. HG-DAgger cedes control to a
human on demand \cite{kelly2019hgdagger}, LazyDAgger cuts context switching
\cite{hoque2021lazydagger}, and ThriftyDAgger budgets interventions by novelty and
risk (Eq.~\eqref{eq:thrifty}) \cite{hoque2021thriftydagger}. DART instead injects
optimized noise into demonstrations for shift-robust offline cloning
\cite{laskey2017dart}, while a deployment-time thread reweights samples around human
takeovers: IWR up-weights bottleneck states \cite{mandlekar2020iwr} and Sirius uses
an approximated trust signal \cite{liu2023sirius}; the landscape is surveyed by
\citet{celemin2022iil}. We compare against the robot-gated subset sharing our
one-demo-per-round budget; human-gated (HG-/Lazy-DAgger), noise-injection (DART),
and reweighting (IWR, Sirius) methods assume a different budget and are discussed
but not run.

Every method above picks one step
$t^\star\!=\!\min\{t:\mathrm{Query}(s_t)\!=\!1\}$ and collects one on-policy segment
per round (Eq.~\eqref{eq:iil-loop}). PACE departs on two counts plus one option: it
localizes failures over completed rollouts and (for prescription) describes them
with a VLM (Eq.~\eqref{eq:perceive}); its central addition clusters failures into
modes and spends the query on a memory-rotated mode
(Eqs.~\eqref{eq:partition}--\eqref{eq:diversity}); and it optionally prescribes the
correction's \emph{start state} rather than resuming at $t^\star$
(Eqs.~\eqref{eq:prescribe}--\eqref{eq:prescribe-demo}). A query-based round is the
corner $|\mathcal{S}|\!=\!1$, $C_{\mathrm{tgt}}\!=\!C^\star$ with an identity
prescribe map, so the DAgger family sits \emph{inside} PACE. The learner, a
diffusion policy, is itself one baseline's query source and thus shared
infrastructure: denoising diffusion models \cite{ho2020ddpm} were adapted to
trajectory planning \cite{janner2022diffuser} and to visuomotor control, capturing
multimodal actions unimodal regression collapses \cite{chi2023diffusionpolicy}, with
Implicit BC contributing the PushT task \cite{florence2021implicitbc}.
Dropout/ensemble gates rest on predictive uncertainty
\cite{gal2016dropout,lakshminarayanan2017ensembles,%
menda2017dropoutdagger,menda2019ensembledagger}, brittle for multimodal policies
where mode disagreement is misread as doubt \cite{lee2025diffdagger}; Diff-DAgger
instead thresholds the per-step diffusion loss against a training-loss quantile
(Eqs.~\eqref{eq:diffsignal}, \eqref{eq:diffdagger}) \cite{lee2025diffdagger}. PACE
reuses this loss $\ell_t$ only to locate each failure's peak frame
(Eq.~\eqref{eq:perceive}) and keys on failure \emph{outcomes} and geometry
(Eq.~\eqref{eq:phi}): it \emph{groups} failures instead of \emph{triggering} on
per-state uncertainty.

\subsection{Foundation Models for Robot Perception and Prescription}
Foundation models increasingly supply robot behavior. SayCan grounds LLM plans in
affordances \cite{ahn2022saycan}, Inner Monologue closes a textual-feedback loop
\cite{huang2022innermonologue}, and PaLM-E and RT-2 fuse perception with language
\cite{driess2023palme,brohan2023rt2}. A separate line emits \emph{executable
structure} --- policy code
\cite{liang2023codeaspolicies,singh2023progprompt}, 3D value maps
\cite{huang2023voxposer}, or reward code
\cite{ma2024eureka,yu2023language2rewards} --- and a recent thread reasons over
failures: REFLECT narrates and corrects robot failures \cite{liu2023reflect}, AHA
reasons over manipulation failures \cite{duan2025aha}, and Reflexion and Self-Refine
iterate on their own language outputs
\cite{shinn2023reflexion,madaan2023selfrefine}. These motivate PACE's use of a
vision--language model (Qwen3-VL \cite{bai2025qwen3vl}) to \emph{describe} why a
rollout failed, feeding only the \emph{optional} prescription stage; unlike
single-trajectory recovery, PACE aggregates the descriptions into a clustered
feature set. When enabled, prescription narrows the LLM's (Qwen3
\cite{yang2025qwen3}, thinking/non-thinking) generativity: rather than plans, code,
or rewards, it emits a compact \texttt{SceneCommand} that the prescribe map
$g(\cdot)$ turns into a single $L_2$-capped, workspace-clamped sub-task-entry start
state (Eq.~\eqref{eq:prescribe}) --- not an executed behavior.

\subsection{Active Learning, Representations, and Reset Curricula}
PACE's \emph{which} decision is batch active learning over failures: pool-based
active learning selects informative examples under a budget \cite{settles2009active},
the core-set/$k$-center formulation maximizes min coverage \cite{sener2018coreset},
BADGE unifies uncertainty and diversity \cite{ash2020badge}, and farthest-point
sampling gives deterministic min-max coverage \cite{eldar1997fps}. PACE clusters
failures into modes with silhouette-chosen $K$-means \cite{lloyd1982kmeans}
(Eq.~\eqref{eq:partition}) and a farthest-point context set (Eq.~\eqref{eq:diversity});
whereas classical methods select \emph{unlabeled inputs}, PACE selects \emph{failure
modes}. We disambiguate two R3M roles: (i) a \emph{frozen}, PCA-reduced embedding
\cite{nair2022r3m} is the visual failure descriptor $\tilde v_i$ that Assess/Choose
cluster over (Eq.~\eqref{eq:phi}), with MVP \cite{radosavovic2022mvp} and VIP
\cite{ma2023vip} as alternatives; (ii) a \emph{separately finetuned} R3M is the
diffusion-policy image encoder (\S Policies per Modality). Our simulators build on
ManiSkill
\cite{mu2021maniskill,gu2023maniskill2,tao2024maniskill3} (Push, from Implicit BC
\cite{florence2021implicitbc}), robomimic \cite{mandlekar2021robomimic}, and
robosuite \cite{zhu2020robosuite} (planned Lift/Wipe/Door). Finally, PACE's optional
Execute stage relates to reset-state and curriculum generation --- reverse
curriculum grows start states outward from the goal
\cite{florensa2017reversecurriculum} and learning-to-reset enables autonomous
operation \cite{eysenbach2018leavenotrace} --- but PACE's entry states follow
clustered failure modes, not a difficulty curriculum; we include a reverse-curriculum
arm as a single-axis \emph{where} control (\S Baselines). Having placed PACE against
each thread, we now define it formally.

\section{Method}
\label{sec:method}

We present \textbf{PACE} (\textbf{P}erceive $\rightarrow$ \textbf{A}ssess
$\rightarrow$ \textbf{C}hoose $\rightarrow$ \textbf{E}xecute), an interactive
imitation-learning (IIL) loop for diffusion policies. \emph{The DAgger-family
baselines answer only \textbf{when} to query the expert; PACE additionally
decides \textbf{which} failure mode to correct --- its central addition --- and,
as an option, \textbf{where} the corrective demonstration should begin.} We first
fix the problem setup and the diffusion-policy learner (Preliminaries), then cast
the classical query rules as instances of one shared decision family (A Unified
Query Framework), and finally define PACE, which extends that family and contains
it as one corner. The four stages map to the equation blocks below: Perceive
(Eqs.~\ref{eq:perceive}--\ref{eq:phi}), Assess
(Eqs.~\ref{eq:partition}--\ref{eq:dominant}), Choose
(Eqs.~\ref{eq:memory}--\ref{eq:diversity}), and Execute
(Eqs.~\ref{eq:prescribe}--\ref{eq:prescribe-demo}).

\subsection{Preliminaries}

\paragraph{Sequential decision process.}
Each task is a finite-horizon (PO)MDP
\begin{equation}
\mathcal{M} = \big(\, \mathcal{S},\ \mathcal{A},\ P(s'\mid s,a),\ R(s,a),\ H,\ \gamma_{\mathrm{env}} \,\big),
\qquad
\pi_\theta:\ \mathcal{S}\to\Delta(\mathcal{A}),
\label{eq:mdp}
\end{equation}
where the state $s$ is privileged proprioception in the state-based setting and a
partial image observation in the image-based setting, $a$ is a discrete move (toy
grid) or a relative joint-position delta (manipulation), and $\pi_\theta$ is the
learner (novice) policy. An expert oracle $\pi^\star$ (A$^\star$/BFS on the grid,
PPO / motion-planner on the robots) induces demonstrations
\begin{equation}
\pi^\star=\arg\max_{\pi}\ \mathbb{E}_{\tau\sim\pi}\!\Big[\textstyle\sum_{t=0}^{H}\gamma_{\mathrm{env}}^{\,t} R(s_t,a_t)\Big],
\qquad
\mathcal{D}=\big\{(s_t,a_t)\ :\ a_t=\pi^\star(s_t)\big\},
\label{eq:expert}
\end{equation}
and behavior cloning fits the policy to the aggregated data,
\begin{equation}
\theta^\star=\arg\min_{\theta}\ \mathbb{E}_{(s,a)\sim\mathcal{D}}\big[\,\mathcal{L}_{\mathrm{BC}}\big(\pi_\theta(\cdot\mid s),\,a\big)\,\big],
\label{eq:bc}
\end{equation}
with $\mathcal{L}_{\mathrm{BC}}=-\log\pi_\theta(a\mid s)$ (discrete) or
$\|\pi_\theta(s)-a\|_2^2$ (continuous). Because the learner visits states its own
errors produce, offline cloning suffers covariate shift
\cite{pomerleau1988alvinn,pomerleau1991efficient}; DAgger corrects this by
aggregating expert labels on \emph{learner-visited} states
\cite{ross2011dagger,ross2014aggrevate,sun2017aggrevated}.

\paragraph{Interactive-IL loop.}
All methods we study share one loop skeleton. In round $r$ the learner
$\pi_{\theta_r}$ is rolled out, a \emph{query predicate} $\mathrm{Query}(s_t)$
flags a first intervention step $t^\star$, the expert takes over from there, and
the resulting expert segment is the single new demonstration:
\begin{align}
\text{roll } \pi_{\theta_r}\ &\Rightarrow\ t^\star=\min\{t:\ \mathrm{Query}(s_t)=1\}, \nonumber\\
\mathcal{D}_{r+1}&=\mathcal{D}_r\ \cup\ \big\{(s_t,\pi^\star(s_t)):\ t\ge t^\star\big\},\quad
\theta_{r+1}=\arg\min_{\theta}\ \mathbb{E}_{\mathcal{D}_{r+1}}[\mathcal{L}].
\label{eq:iil-loop}
\end{align}
Here $\mathcal{L}$ is the imitation loss: $\mathcal{L}_{\mathrm{BC}}$
(Eq.~\ref{eq:bc}) for the T1 classifier/CNN and $\mathcal{L}_{\mathrm{DP}}$
(Eq.~\ref{eq:vpred}) for the T2 (and planned T3--T5) diffusion learner. Exactly one
demonstration is added per round, the policy is retrained from scratch every $n_d$
demos, and the loop stops when a held-out success rate
$\mathrm{SR}_{\mathrm{target}}=0.90$ is reached or a budget $B$ is spent. This
one-demo-per-round protocol is the shared sample-efficiency axis on which every
method is compared, and the only degree of freedom is how the new demonstration is
chosen.

\paragraph{Diffusion-policy learner.}
Every robot-task method (PACE, Diff-DAgger, and all IIL baselines) shares one
diffusion-policy backbone \cite{chi2023diffusionpolicy,ho2020ddpm,janner2022diffuser}
so that only the demo-acquisition rule differs. With clean action target
$x_0=a$, the forward process noises $x_0$ over $K_{\mathrm{dif}}$ steps,
\begin{equation}
x_k=\sqrt{\bar\alpha_k}\,x_0+\sqrt{1-\bar\alpha_k}\,\epsilon,
\qquad \epsilon\sim\mathcal{N}(0,I),\quad k\in\{1,\dots,K_{\mathrm{dif}}\},
\label{eq:forward}
\end{equation}
and the network is trained with a $v$-prediction denoising objective (conditioning
on $s$ suppressed for brevity),
\begin{equation}
v_k=\sqrt{\bar\alpha_k}\,\epsilon-\sqrt{1-\bar\alpha_k}\,x_0,
\qquad
\mathcal{L}_{\mathrm{DP}}(\theta)=\mathbb{E}_{k,\epsilon,(s,a)}\big[\big\|v_\theta(x_k,k,s)-v_k\big\|_2^2\big].
\label{eq:vpred}
\end{equation}
Crucially, the per-pair denoising loss doubles as an uncertainty signal
\cite{lee2025diffdagger}: averaged over noise draws it scores how out-of-distribution
a state--action pair is, and evaluated at the novice's own executed action along a
rollout it gives the per-step signal $\ell_t$,
\begin{equation}
L_{\mathrm{dif}}(s,a)=\mathbb{E}_{k,\epsilon}\big[\big\|v_\theta(x_k,k,s)-v_k\big\|_2^2\big],
\qquad
\ell_t \;=\; L_{\mathrm{dif}}\big(s_t,\,a^{\mathrm{nov}}_t\big).
\label{eq:diffsignal}
\end{equation}
$\ell_t$ is the trigger for Diff-DAgger and, in PACE, the analysis frame selector
that localizes each failure (Perceive).

\subsection{A Unified Query Framework for the Baselines}
\label{sec:unified}

Every IIL baseline instantiates the same template: it computes a scalar score at
each visited state and queries when that score crosses a threshold,
$\mathrm{Query}_{\mathrm{method}}(s_t)=\mathbf{1}\big[\mathrm{score}_{\mathrm{method}}(s_t)\gtrless\tau_{\mathrm{method}}\big]$.
All action comparisons are made in the shared $9$-D relative-joint-position space. The
PPO expert's $7$ arm-joint deltas are lifted to $9$-D by a fixed, method-independent
embedding whose two extra \texttt{panda\_stick} finger dimensions take the novice's
own values (zero contribution to $\|a^{\mathrm{nov}}_t-a^{\mathrm{exp}}_t\|_2$); the
same lift is used by all discrepancy-based methods, so none is advantaged. The score
distinguishes the methods: \emph{SafeDAgger}$^\star$ on action discrepancy
(Eq.~\ref{eq:safe}); \emph{DropoutDAgger} on MC-sampled in-ball agreement
(Eq.~\ref{eq:dropout}); \emph{EnsembleDAgger} on ensemble doubt or mean discrepancy
(Eq.~\ref{eq:ensemble}); \emph{ThriftyDAgger} on novelty or learned task-risk $1-Q$
(Eq.~\ref{eq:thrifty}); and a uniform-random \emph{Random} control
(Eq.~\ref{eq:stagger}). We defer these five predicates and their constants to the
Appendix and keep the diffusion-native rule in the body, since it is the one PACE
reduces toward. \emph{Diff-DAgger} \cite{lee2025diffdagger} uses the diffusion loss
of Eq.~\ref{eq:diffsignal}: with $\eta_\alpha$ the $\alpha$-quantile of the
training-loss CDF (recalibrated each retrain), it queries when the last $K$ steps are
all in the tail,
\begin{equation}
\eta_\alpha=\hat F^{-1}(\alpha),
\qquad
\mathrm{Query}_{\mathrm{Diff}}(s_t)=\mathbf{1}\!\left[\ \sum_{u=t-K+1}^{t}\mathbf{1}\big[\ \ell_u>\eta_\alpha\ \big]\ \ge\ K\ \right].
\label{eq:diffdagger}
\end{equation}
Two faithfulness notes. \emph{SafeDAgger}$^\star$ is starred because we gate on the
raw novice--expert action discrepancy against the (A$^\star$/planner) oracle rather
than the learned safety classifier of the original~\cite{zhang2017safedagger}: same
decision variable, an oracle signal in place of a learned one; we do not claim this
is conservative and label it a variant of unknown bias. For \emph{ThriftyDAgger}
\cite{hoque2021thriftydagger}, the success-to-go discount $\gamma_Q$ is distinct from
the MDP discount $\gamma_{\mathrm{env}}$ (Eq.~\ref{eq:mdp}) and the memory discount
$\gamma_{\mathrm{mem}}$ (Eq.~\ref{eq:memory}). Both baselines are fully specified in
the Appendix.

\paragraph{What the framework leaves on the table.}
Eqs.~\ref{eq:safe}--\ref{eq:stagger} and \ref{eq:diffdagger} answer only \emph{at
which step $t^\star$ should the expert intervene?} and then correct on-policy from
that step. Two limitations follow, and they are what PACE's \emph{which} axis
targets: (i) being per-state, they cannot reason about which failures in a batch are
redundant, and (ii) being memoryless across rounds, they may correct a persistent
mode repeatedly. A softer third gap is that they inherit the start state of the
tripping rollout; PACE can optionally place the demo at a prescribed sub-task entry
point (the \emph{where} option), which the ablations show is not the driver of any
gain. PACE keeps the loop of Eq.~\ref{eq:iil-loop} unchanged and replaces only the
single-scalar query with a four-stage pipeline; the next subsection defines each
stage and Algorithm~\ref{alg:pace} assembles them.

\subsection{PACE}
\label{sec:pace}

PACE operates on the \emph{set} of failures $\mathcal{F}_r=\{f_1,\dots,f_N\}$
produced by a round's rollouts, transforming it into one corrective demonstration;
Algorithm~\ref{alg:pace} gives the full loop.

\paragraph{Perceive.}
For each failure $i$ we localize the failure point $t^\star_i$ as the peak-loss step
of Eq.~\ref{eq:diffsignal} (maximal diffusion uncertainty) and record raw geometry
there:
\begin{equation}
\mathcal{F}_r=\{f_1,\dots,f_N\},
\quad
t^\star_i=\arg\max_t \ell^{\,(i)}_t,\quad
\mathrm{peak}_i=\max_t \ell^{\,(i)}_t,\quad
\rho_i=\frac{t^\star_i}{\max(1,T_i)},\quad
\delta_i=\big\|p^{\mathrm{tcp}}_i-p^{\mathrm{tee}}_i\big\|_2 .
\label{eq:perceive}
\end{equation}
Here $\rho_i$ is task progress and $\delta_i$ a contact distance between object and
end-effector. Equations \ref{eq:perceive} and \ref{eq:prescribe} (the geometry-bearing
Perceive and Execute terms) are written for the \emph{manipulation} instantiation
(T2, and T3--T5 when wired), where
$p^{\mathrm{tee}}_i$ is the object's planar pose and $p^{\mathrm{tcp}}_i$ the
end-effector (tool-center-point) pose at $t^\star_i$. The toy grid (T1) uses a
distinct instantiation of the same pipeline (grid-cell geometry, classifier-BCE
$\ell_t$, corridor prescription) with the Assess/Choose stages
(Eqs.~\ref{eq:partition}--\ref{eq:diversity}) unchanged; the full grid walkthrough is
in the Appendix. A
vision--language model \cite{bai2025qwen3vl,duan2025aha,liu2023reflect} reads three
key frames around $t^\star_i$ (start, peak, end) with a per-task knowledge-augmented
generation (KAG) context --- a knowledge graph of workspace bounds and a failure
taxonomy --- and \emph{describes} why the rollout failed; this text is the VLM's sole
output into the loop. \textbf{Where the VLM acts, precisely.} Localization uses the
scalar peak-loss step (Eq.~\ref{eq:diffsignal}) and mode discovery (Assess, Choose)
runs on the descriptor $\psi_i$ below, both without a VLM; the VLM text enters only
Execute, where it conditions the LLM \texttt{SceneCommand} that $g(\cdot)$ turns into
the corrective reset (Eq.~\ref{eq:prescribe})
\cite{driess2023palme,huang2022innermonologue,shinn2023reflexion,madaan2023selfrefine},
and a dedicated VLM-off ablation (\S Ablations) isolates whether it helps. Failures
are featurized by a quaternion-free geometric descriptor $\phi_i\in\mathbb{R}^6$,
with an optional \emph{frozen} R3M visual embedding \cite{nair2022r3m} (PCA-reduced)
for the image-based setting --- distinct from the separately finetuned R3M encoder
used inside the diffusion policy (\S Policies per Modality):
\begin{align}
\phi_i&=\big[\,p^{\mathrm{tee}}_{i,x},\ p^{\mathrm{tee}}_{i,y},\ \sin\theta_i,\ \cos\theta_i,\ \rho_i,\ \delta_i\,\big]\in\mathbb{R}^6,
\quad
\tilde v_i=\mathrm{PCA}_{k'}\big(v_i\big),\ \ k'=\min(16,N-1,d), \nonumber\\
\psi_i&=\begin{cases}\tilde v_i & \text{visual mode}\\[2pt]\phi_i & \text{geometric (default)}\end{cases},
\quad
X=[\psi_1;\dots;\psi_N],\quad
\tilde X_i=\big(\psi_i-\mu\big)\oslash\sigma,\quad \sigma_j\!\leftarrow\!\max(\sigma_j,10^{-8}).
\label{eq:phi}
\end{align}
The feature $\psi_i$ is used only for mode discovery; all prescription geometry
(Execute) uses the raw object pose and full robot configuration.

\paragraph{Assess.}
PACE clusters the standardized features into failure modes, choosing the number of
clusters by maximum mean silhouette \cite{lloyd1982kmeans} and computing per-cluster
geometry on the raw object pose:
\begin{align}
k^\star&=\arg\max_{k\in[2,\,k_{\max}]}\ \mathrm{sil}(k),\qquad k_{\max}=\max\!\big(2,\min(6,N-1)\big),\nonumber\\
c^{xy}_{C}&=\tfrac{1}{|C|}\!\sum_{i\in C}p^{\mathrm{tee}}_{i,xy},\quad
\bar\theta_{C}=\operatorname{atan2}\!\Big(\!\sum_{i\in C}\!\sin\theta_i,\ \sum_{i\in C}\!\cos\theta_i\Big),\quad
\bar L_{C}=\tfrac{1}{|C|}\!\sum_{i\in C}\mathrm{peak}_i,\nonumber\\
\mathrm{rep}(C)&=\arg\min_{i\in C}\ \big\|\tilde X_i-\bar X_{C}\big\|_2,\qquad
\bar X_{C}=\tfrac1{|C|}\!\sum_{i\in C}\tilde X_i.
\label{eq:partition}
\end{align}
The cluster representative $\mathrm{rep}(C)$ is the member nearest the cluster mean
in feature space. The \emph{dominant} mode --- the most prevalent, hardest failure
mode --- is selected by a deterministic lexicographic key on (size, mean peak-loss,
earliest episode),
\begin{equation}
C^\star=\operatorname*{arg\,lexmax}_{C\in\{C_1..C_m\}}\ \big(\ |C|,\ \bar L_{C},\ -\!\min_{i\in C}\mathrm{eid}_i\ \big).
\label{eq:dominant}
\end{equation}

\paragraph{Choose.}
This stage decides \emph{which} failure mode the round corrects. Because the policy
is retrained after every demo, the failure distribution shifts each round and a
single mode could be prescribed repeatedly. PACE therefore keeps a cross-round
\emph{memory} of already-corrected centroids and applies a recency-discounted
Gaussian coverage penalty, rotating the target mode among those tied with the
dominant one to avoid re-covering a solved mode
\cite{florensa2017reversecurriculum,eysenbach2018leavenotrace}:
\begin{align}
P_{\mathrm{mem}}(c)&=\sum_{i:\,(r_i,c_i)\in\mathrm{Mem}}\gamma_{\mathrm{mem}}^{\,\max(0,\,r-r_i)}\exp\!\Big(-\tfrac{\|c-c_i\|_2^2}{2\sigma_{\mathrm{mem}}^2}\Big),\nonumber\\
C_{\mathrm{tgt}}&=\arg\max_{C:\,|C|\ge|C^\star|-1}\ \Big(\ \bar L_{C}\ -\ \lambda\,P_{\mathrm{mem}}\big(c^{xy}_{C}\big)\ \Big),\qquad \gamma_{\mathrm{mem}}=0.6,\ \sigma_{\mathrm{mem}}=0.06,\ \lambda=1.
\label{eq:memory}
\end{align}
Given the target, PACE selects the within-round set by a memoryless \emph{coreset /
k-center} (farthest-point) rule
\cite{sener2018coreset,eldar1997fps,ash2020badge,settles2009active}, forcing the
target representative in and seeding the hardest failure:
\begin{align}
S_0&=\big\{\,\mathrm{rep}(C_{\mathrm{tgt}})\,\big\}\ \cup\ \big\{\,\arg\max_i \mathrm{peak}_i\,\big\}, \nonumber\\
S&\leftarrow S\ \cup\ \Big\{\ \arg\max_{i\notin S}\ \min_{j\in S}\ \big\|\tilde X_i-\tilde X_j\big\|_2\ \Big\}\quad\text{until}\ |S|=\kappa,\qquad \kappa=3, \nonumber\\
S^\star&=\arg\max_{\substack{S\subseteq\mathcal{F}_r,\ |S|\le\kappa\\ S_0\subseteq S}}\ \min_{\substack{i,j\in S\\ i\ne j}}\ \big\|\tilde X_i-\tilde X_j\big\|_2
\quad\text{s.t.}\ \ S_0=\big\{\mathrm{rep}(C_{\mathrm{tgt}}),\ \arg\max_i\mathrm{peak}_i\big\}.
\label{eq:diversity}
\end{align}
The last line is the max-min coverage objective PACE approximates; the first two are
its greedy realization, and forcing the same seed $S_0$ (target representative and
worst-loss failure) into both makes the greedy pick optimize exactly the constrained
objective. Under a one-demo-per-round budget the coverage work is done by the memory
rotation of Eq.~\ref{eq:memory}, which changes \emph{which mode} $C_{\mathrm{tgt}}$
anchors each round and accumulates a broad, non-redundant cover across rounds. The
within-round set $S$ of Eq.~\ref{eq:diversity} is subordinate:
$\mathrm{rep}(C_{\mathrm{tgt}})$ is the demo anchor and the remaining $\kappa{-}1$
members are cited context and escalation fallback. Under the default on-policy
Execute they do not change the collected demo (a tie-break); they matter only under
a \texttt{BRIDGE} prescription, where the LLM places a middle-ground start near
$2$--$3$ cited members so the context set alters \emph{where} the demo begins. This
is the cross-round coverage axis the per-state baselines of \S\ref{sec:unified} lack.
The ablation separates these (Table~\ref{tab:ablation}): a $-$Memory row isolates the
acting mechanism and a $-$Choose-context row plus the BRIDGE-on/off contrast tests
whether the context set ever changes the demo. Choose fixes \emph{which} mode anchors
the round; Execute decides \emph{where} the corrective demonstration begins.

\paragraph{Execute.}
The anchor $A$ is the target representative's raw geometry. Execute collects the one
new demonstration and has a default and an option. \textbf{Default (on-policy,
$g=\mathrm{identity}$).} The demonstration re-rolls the anchor failure's own scene
and lets the expert take over in place from step $t^\star_A$ --- exactly the
DAgger-family on-policy correction, with the \emph{which} decision applied on top.
This keeps the \emph{where} decision off, so PACE's gain must come entirely from
\emph{which} mode is corrected; we treat it as the primary configuration.
\textbf{Option (prescribed \emph{where}, $g\neq\mathrm{identity}$).} When enabled, a
VLM/LLM
\cite{yang2025qwen3,ma2024eureka,yu2023language2rewards,liang2023codeaspolicies}
emits a scene command $\mathrm{cmd}$ (target object pose, TCP, decision label) that
$g(\cdot)$ turns into a sub-task-entry start state $\xi$. When $\mathrm{cmd}$ points
at a single recorded failure the round reduces to the default (a \texttt{SELECT});
when it proposes a middle-ground pose over $2$--$3$ cited members it becomes a
\texttt{BRIDGE}, the only mode that genuinely uses $g\neq\mathrm{identity}$. The
anchor-relative perturbation is $L_2$-capped in position and clamped in yaw, and the
spec is projected onto the KAG workspace bounds $\mathcal{W}$:
\begin{align}
A&=\big(p^{\mathrm{tee}}_A,\ \theta_A,\ q^{\mathrm{full}}_A,\ p^{\mathrm{tcp}}_A,\ t^\star_A\big)=\mathrm{rep}(C_{\mathrm{tgt}}),
\quad
\mathrm{cmd}=\mathrm{VLM/LLM}\big(A,\ \text{frames}(t^\star_A),\ \mathrm{KAG}\big),\nonumber\\
\Delta_{xy}&=\mathrm{cmd}.p^{\mathrm{tee}}_{xy}-p^{\mathrm{tee}}_{A,xy},\quad
\Delta_{xy}\leftarrow\Delta_{xy}\cdot\min\!\Big(1,\ \tfrac{\Delta_{\max}}{\|\Delta_{xy}\|_2}\Big),\quad
\Delta_\theta=\mathrm{clamp}\big(\mathrm{wrap}_\pi(\mathrm{cmd}.\theta-\theta_A),\,\pm\theta_{\max}\big),\nonumber\\
\xi&=g(\mathrm{cmd})=\mathrm{clamp}_{\,\mathcal{W}}\Big(p^{\mathrm{tee}}_A+[\Delta_{xy},0],\ \ \theta_A+\Delta_\theta,\ \ q^{\mathrm{full}}_A\Big),
\quad
(\Delta_{\max},\theta_{\max})=(0.06,0.4).
\label{eq:prescribe}
\end{align}
Either way, the round samples the start state from the reset distribution
$\Xi(\cdot\mid\xi)$ ($\xi$ is the on-policy anchor state under the default), lets
the expert roll from $t^\star$, and aggregates the single new demo:
\begin{equation}
s_0\sim\Xi(\,\cdot\mid\xi\,),
\qquad
\mathcal{D}_{r+1}=\mathcal{D}_r\ \cup\ \big\{(s_t,\pi^\star(s_t)):\ t\ge t^\star,\ s_0\sim\Xi(\cdot\mid\xi)\big\}.
\label{eq:prescribe-demo}
\end{equation}
An empty or refused prescription is never accepted (a hard non-emptiness floor with
bounded re-prescription attempts), since a wasted round is the worst outcome under a
fixed budget. The caps $(\Delta_{\max},\theta_{\max})$ and memory constants
$(\gamma_{\mathrm{mem}},\sigma_{\mathrm{mem}},\lambda)$ are given for Push; per-task
values are carried in the configuration. Because the default already carries the
\emph{which} decision, the $-$Execute ablation (Table~\ref{tab:ablation}) is exactly
the test of whether the optional \emph{where} adds anything on top; if not, the
contribution rests on \emph{which}.

\paragraph{A unifying view: the baselines as a corner of PACE.}
We close with an organizing observation (not a contribution), scoped to the
diffusion-loss task (T2): the query-based family sits at one corner of the PACE
design space. A PACE round collapses to a single on-policy query under three
\emph{simultaneous} restrictions: (i) $|S|=1$; (ii) memory and rotation off,
$C_{\mathrm{tgt}}=C^\star$ reduced to the peak-loss failure; and (iii) the on-policy
default $g=\mathrm{identity}$, so $\Xi(\cdot\mid\xi)$ replays the failing rollout and
the expert corrects in place from $t^\star$. Requiring all three at once, this is a
corner of a three-parameter family, not a single-axis nesting; PACE contains the
query family at this corner and we claim nothing stronger. Even here, equality is
delicate because PACE's discovery differs structurally: it analyzes \emph{completed}
rollouts and localizes the peak-loss frame $t^\star_i=\arg\max_t\ell^{(i)}_t$
(Eq.~\ref{eq:perceive}), whereas Diff-DAgger triggers \emph{online} at the first
windowed tail-crossing (Eq.~\ref{eq:diffdagger}). The relation is a definitional
inequality, not a result: with window $K{=}1$ and a non-empty tail, the first
crossing $\min\{t:\ell_t>\eta_\alpha\}$ never follows the maximum $\arg\max_t\ell_t$,
and the two coincide exactly when the loss is non-decreasing up to its peak (then the
rounds aggregate the identical demonstration; otherwise their $t^\star$ differ). It
does not extend to the other baselines, whose triggers use different scores
(Eqs.~\ref{eq:safe}--\ref{eq:stagger}) and are recovered by the same corner only up
to substituting their own $\mathrm{Query}(\cdot)$, and is vacuous on the toy grid
(Appendix). What PACE \emph{adds} over the corner is cross-round memory
rotation (Eq.~\ref{eq:memory}), the within-round context set (Eq.~\ref{eq:diversity},
active only under \texttt{BRIDGE}), and the optional learned \emph{where}-to-start
prescription ($g\neq\mathrm{identity}$, Eq.~\ref{eq:prescribe}) --- with the VLM/LLM
grounding, all ablated in \S Ablations.

\begin{algorithm}[tb]
\caption{PACE: Perceive $\rightarrow$ Assess $\rightarrow$ Choose $\rightarrow$
Execute. Written modality-agnostically ($\mathcal{L}$ is the task loss); the T1
grid replaces the Perceive signal and the prescribe map $g$ as described in the
Method, keeping the Assess/Choose core.}
\label{alg:pace}
\textbf{Input}: initial demos $\mathcal{D}_0$, expert $\pi^\star$, budget $B$,
retrain cadence $n_d$, target $\mathrm{SR}_{\mathrm{target}}$, cap $\kappa$\\
\textbf{Output}: trained policy $\pi_\theta$
\begin{algorithmic}[1]
\STATE $\theta \leftarrow \arg\min_\theta \mathbb{E}_{\mathcal{D}_0}[\mathcal{L}]$;\ \ $r\leftarrow 0$;\ \ $q\leftarrow 0$;\ \ $\mathrm{Mem}\leftarrow\varnothing$ \hfill{\footnotesize// $\mathcal{L}{=}\mathcal{L}_{\mathrm{DP}}$ (robot, Eq.~\ref{eq:vpred}) or $\mathcal{L}_{\mathrm{BC}}$ (T1)}
\WHILE{$q<B$ \textbf{and} $\mathrm{SR}(\theta)<\mathrm{SR}_{\mathrm{target}}$}
  \STATE roll $\pi_\theta$; collect failure set $\mathcal{F}_r$; for each $f_i$ get $t^\star_i,\mathrm{peak}_i,\rho_i,\delta_i$ \hfill{\footnotesize// \textsc{Perceive}, Eq.~\ref{eq:perceive}}
  \STATE VLM describes failures around $t^\star_i$; featurize $\psi_i$; standardize to $\tilde X$ \hfill{\footnotesize// Eq.~\ref{eq:phi}}
  \STATE cluster $\{C_1,\dots\}$ at $k^\star$; compute $\mathrm{rep}(C),\bar L_C$; pick dominant $C^\star$ \hfill{\footnotesize// \textsc{Assess}, Eqs.~\ref{eq:partition}--\ref{eq:dominant}}
  \STATE $C_{\mathrm{tgt}}\leftarrow$ memory-rotated target (coverage mechanism);\ \ $S\leftarrow$ k-center context, $\mathrm{rep}(C_{\mathrm{tgt}})$ forced, $|S|\!\le\!\kappa$ \hfill{\footnotesize// \textsc{Choose}, Eqs.~\ref{eq:memory}--\ref{eq:diversity}}
  \STATE \textbf{if} where-option off: $\xi\leftarrow$ anchor's own $t^\star_A$ state ($g{=}\mathrm{id}$); \textbf{else} $\mathrm{cmd}\leftarrow$ VLM/LLM$(A,\mathrm{frames},\mathrm{KAG})$, $\xi\leftarrow g(\mathrm{cmd})$ \hfill{\footnotesize// \textsc{Execute}, Eq.~\ref{eq:prescribe}}
  \STATE $s_0\!\sim\!\Xi(\cdot|\xi)$; expert rolls from $t^\star$; $\mathcal{D}_{r+1}\!\leftarrow\!\mathcal{D}_r\cup\{(s_t,\pi^\star(s_t)):t\ge t^\star\}$ \hfill{\footnotesize// Eq.~\ref{eq:prescribe-demo}}
  \STATE append $(r,c^{xy}_{C_{\mathrm{tgt}}},\bar\theta_{C_{\mathrm{tgt}}})$ to $\mathrm{Mem}$;\ \ $q\leftarrow q+1$;\ \ $r\leftarrow r+1$
  \IF{$q \bmod n_d = 0$}
    \STATE $\theta \leftarrow \arg\min_\theta \mathbb{E}_{\mathcal{D}_r}[\mathcal{L}]$ \hfill{\footnotesize// retrain from scratch; recalibrate $\eta_\alpha$ (robot)}
  \ENDIF
\ENDWHILE
\STATE \textbf{return} $\pi_\theta$
\end{algorithmic}
\end{algorithm}

\section{Experiments}

The unifying view of \S\ref{sec:method} anchors the T2 comparison: at its corner PACE
reduces to Diff-DAgger, so a measured gain there isolates the added decision. That
corner is built on the diffusion loss and applies to T2 only; Diff-DAgger is absent
from the toy grid (T1), which shares only the Assess/Choose core under a distinct
Perceive signal (classifier BCE loss) and Execute map. T1 is thus a weaker sanity
check on whether clustering$+$memory-rotation helps a discrete IIL task, read against
the diffusion-free analogue ``correct the highest-loss failure on-policy'' (the
$-$Assess/$-$Choose/$-$Execute reduction). We test PACE along four axes.
\textbf{(Q1, primary)}~Under a fixed one-demonstration-per-round budget, does PACE
reach a target held-out success rate with fewer expert queries than uncertainty- and
safety-gated IIL baselines and the IWR-style arm? \textbf{(Q2, secondary/optional)}~Does
the optional prescribed start-state (\texttt{BRIDGE}) improve efficiency and coverage
on top of the on-policy default, or does the reset-curriculum arm already capture that
gain? We treat Q2 as a test of an option, not the thesis. \textbf{(Q3)}~Do the Q1
gains hold across observation modalities (privileged state vs.\ raw image) on the two
implemented tasks (T1 grid, T2 Franka Panda)? Cross-embodiment transfer to UR5/UR5e is
left to the planned RoboSuite extension (\S Task Suite) and is \emph{not} claimed here.
\textbf{(Q4)}~Which of the four PACE stages---Perceive, Assess, Choose, Execute---and
which foundation model (VLM vs.\ LLM) drive any improvement? Because the five-seed
runs are in progress, every quantitative entry is a red placeholder macro and we
phrase each empirical statement as a question the pending numbers will answer; the
tables and callouts fix the reporting protocol (\S Metrics) so no result ships
un-filled or un-tested.

\subsection{Task Suite}

We evaluate on the two implemented tasks summarized in Table~\ref{tab:tasks} (T1
Toy, T2 Push), spanning a discrete navigation problem with an exact expert and a
contact-rich pushing task; three RoboSuite tasks on UR5/UR5e arms (T3--T5) are
specified below as a \emph{planned} protocol only. Both implemented tasks are
genuinely multi-modal in optimal behavior --- the regime in which per-state
action-agreement uncertainty is least reliable and failure-mode partitioning (Assess)
is expected to help.

\textbf{T1 (Toy).} A custom $5{\times}5$ grid-navigation environment
(\texttt{MazeNavEnv}) in which a point agent moves \{up, down, left, right\}
($\mathcal{A}=\mathrm{Discrete}(4)$) from a fixed start to a fixed goal; two central
fire tiles leave two equally-optimal $8$-step routes (top, bottom), making the task
multi-modal by construction. The $A^{\!\star}$/BFS expert emits a multi-label
optimal-action mask (every action that reduces BFS distance-to-goal is positive),
teaching both routes. The state observation is a $14$-d vector (agent/goal
coordinates, a $3{\times}3$ local tile patch, steps-remaining fraction); the image
observation is an $80{\times}80{\times}3$ bird's-eye render. Held-out evaluation uses
$\PH{toy-heldout-n}$ fixed layouts.

\textbf{T2 (Push).} ManiSkill3~\cite{tao2024maniskill3,mu2021maniskill} PushT with a
Franka Panda and a \texttt{panda\_stick} end-effector: push a T-shaped block onto a
fixed goal pose (success $=$ overlap within tolerance); PushT originates in Implicit
BC~\cite{florence2021implicitbc}. State mode exposes the privileged T-pose
(\texttt{proprio\_dim}\,$=21$); image mode uses a $256{\times}256$ base-camera RGB
stream and hides the T-pose (\texttt{proprio\_dim}\,$=14$), forcing the policy to read
the object from pixels. The action space is
$\texttt{rel\_joint\_pos}\in\mathbb{R}^{9}$ and the expert is a privileged PPO policy,
identical across modalities so demonstrations differ only by observation.

\textbf{T3--T5 (Lift, Wipe, Door) --- planned, not evaluated.}
We \emph{specify} but do not run three RoboSuite~\cite{zhu2020robosuite}
manipulation tasks on a UR5e (Lift) and UR5 (Wipe, Door), each intended to use state
and image diffusion policies and a motion-planner expert, following the offline
protocol of robomimic~\cite{mandlekar2021robomimic} and reusing the T2 loop verbatim.
These tasks are \emph{not implemented} (simulator wiring, policy configs, and experts
are pending); we describe them only to fix the intended cross-embodiment protocol.
They contribute \emph{no result rows} and no UR5/UR5e generality claim is made.

\subsection{Policies per Modality}

To isolate the demonstration-acquisition rule, all arms share one learner per
(task, modality) cell. On T1 the \emph{state} learner is a $p4m$/$D_4$
group-equivariant policy (\texttt{escnn}) and the \emph{image} learner a plain
grid-size-invariant CNN. On T2 (and planned T3--T5) both modalities use a single
shared diffusion-policy backbone~\cite{chi2023diffusionpolicy,ho2020ddpm}: a
conditional $1$-D U-Net denoiser (obs horizon $1$, prediction horizon $32$, action
horizon $2$), trained with the $v$-prediction objective (Eq.~\ref{eq:vpred}) under a
DDIM schedule. The state encoder is an MLP; the image encoder is an R3M
ResNet-18~\cite{nair2022r3m} \emph{finetuned end-to-end} with a spatial-softmax head
(output dim $1038$ for PushT) --- a distinct module from the frozen R3M failure
descriptor of Assess/Choose (\S Method). One identical learner per arm keeps the
comparison apples-to-apples in the regime Diffusion Policy and Diff-DAgger
target~\cite{chi2023diffusionpolicy,lee2025diffdagger}.

\subsection{Baselines}

We compare PACE against the interactive-IL family, all cast as a single decision-rule
schema $\mathrm{Query}(s_t)=\mathbf{1}[\,\mathrm{score}_\mathrm{method}(s_t)\ \text{crosses threshold}\,]$
so that only the query criterion changes. \textbf{SafeDAgger}~\cite{zhang2017safedagger}
queries on novice--expert action discrepancy (Eq.~\ref{eq:safe});
\textbf{DropoutDAgger}~\cite{menda2017dropoutdagger,gal2016dropout} fires when
MC-sampled in-ball agreement drops (Eq.~\ref{eq:dropout}); and
\textbf{EnsembleDAgger}~\cite{menda2019ensembledagger,lakshminarayanan2017ensembles}
fires on ensemble doubt or mean discrepancy (Eq.~\ref{eq:ensemble}).
\textbf{ThriftyDAgger}~\cite{hoque2021thriftydagger} queries on novelty or learned
task risk $1-Q_\psi$ with budget-calibrated thresholds (Eq.~\ref{eq:thrifty}).
\textbf{Random} is a uniform one-demo-per-round control (``Stagger'',
Eq.~\ref{eq:stagger}), a floor rather than a published method. On robot tasks we
add \textbf{Diff-DAgger}~\cite{lee2025diffdagger}, the diffusion-native rule that
queries when the per-step diffusion loss stays in the tail of its training-loss CDF
for $K$ windowed steps (Eq.~\ref{eq:diffdagger}); it is omitted on T1, which lacks
the diffusion learner. \textbf{Two single-axis controls} isolate the integration from its parts, since
integration-only novelty is credible only when the integrated system beats the best
individual axis. \textbf{Reset-Curriculum} ($R$-arm) is a reverse-curriculum
control~\cite{florensa2017reversecurriculum}: it grows the corrective start state
outward from the failure state (a Brownian step under the same $L_2$/workspace caps as
Eq.~\ref{eq:prescribe}) \emph{without} the Perceive/Assess/Choose curation, and shares
PACE's reset access. \textbf{IWR-Reweight} ($W$-arm, in the spirit of
IWR~\cite{mandlekar2020iwr}) up-weights failure-region samples in the training loss
\emph{after} on-policy collection at $t^\star$, and uses no reset access. Both run
under the identical one-demo-per-round loop and shared bootstrap, giving T1 seven arms
(five DAgger-family $+$ two controls) and T2 eight ($+$ Diff-DAgger). Each baseline
uses its default hyperparameters, reported \emph{per code path} where the two differ:
SafeDAgger $\tau=0.1$; EnsembleDAgger $M=5$; ThriftyDAgger $\alpha_h=0.1$;
Diff-DAgger $\alpha=0.99$, $K=1$; and \textbf{DropoutDAgger} on robot ($N=10$
diffusion draws, $\tau=0.1$, in-ball fraction threshold $p=0.9$) vs.\ grid ($N=16$
MC-dropout samples, dropout rate $0.2$, agreement threshold $p_{\mathrm{thr}}=0.5$),
the split reflecting that the diffusion policy is natively stochastic while the grid
classifier needs injected dropout. PACE reuses the same per-step signal to
\emph{perceive} failures but also partitions them, selects a memory-rotated mode, and
optionally prescribes the corrective entry state, recovering Diff-DAgger at its corner
($|S|=1$, $C_\mathrm{tgt}=C^\star$, on-policy identity prescribe map,
Eq.~\ref{eq:prescribe-demo}).

\subsection{Active-Loop Protocol}

Every method follows the same interactive loop (Eq.~\ref{eq:iil-loop}): bootstrap a
policy from a small shared set of expert demonstrations, then repeat \{roll out
$\to$ query rule flags a failure $\to$ collect exactly \textbf{one} new expert
demonstration $\to$ retrain $\to$ held-out eval\} until a stop condition. The
\textbf{one-demo-per-round} cap is a hard invariant that makes the number of expert
queries the shared sample-efficiency axis, and all arms share the same bootstrap
policy per seed, so the only difference between arms is the acquisition rule.

\textbf{Toy (T1).} We bootstrap from $\PH{toy-init-demos}$ $A^{\!\star}$/BFS
demonstrations, draw a size-$\PH{toy-pool-n}$ correction pool per round, retrain from
scratch each round, and evaluate on $\PH{toy-heldout-n}$ fixed held-out layouts. The
budget is $\PH{toy-budget}$ extra demonstrations over the bootstrap, with a
$\PH{toy-max-rounds}$-round cap and rollout cap of $\PH{toy-max-steps}$ steps and
target $\mathrm{SR}_\mathrm{target}=0.90$ (Eq.~\ref{eq:sr}). A run stops at target,
budget exhaustion, or round cap.

\textbf{Robot (T2; identical for planned T3--T5).} We bootstrap a behavior-cloned
diffusion policy from $\PH{robot-init-demos}$ expert demonstrations (one shared warm
start per seed). Each round discovers a failure via the method's rule over a
$\PH{robot-rollout-n}$-episode rollout and adds one demonstration; the policy is
retrained \emph{from scratch every} $n_d=\PH{robot-nd}$ \emph{demonstrations}
following Diff-DAgger~\cite{lee2025diffdagger}, a cadence shared across methods. An
episode counts as a demonstration only if the expert succeeded, actually intervened,
and contributed at least $\PH{robot-min-expert-steps}$ expert steps. We evaluate on
$\PH{robot-heldout-n}$ held-out episodes per checkpoint (fixed held-out seed base
$\PH{robot-heldout-seed}$, vectorized over $\PH{robot-eval-envs}$ environments), with
a budget of $\PH{robot-budget}$ additional demonstrations, a $\PH{robot-max-rounds}$-round
cap, and $\mathrm{SR}_\mathrm{target}=0.90$; the query count is recorded when the
target is first reached.

\subsection{Metrics and Seeds}

For each (task, modality, method, seed) we report four metrics from the learning
curve. \textbf{SR}: final held-out success rate (Eq.~\ref{eq:sr}). \textbf{Q}: expert
queries---equivalently demonstrations added---to first reach
$\mathrm{SR}_\mathrm{target}=0.90$ (Eq.~\ref{eq:queries}). $Q$ is
\emph{right-censored}: a run that never reaches target contributes $Q=B$, so a bare
mean$\pm$std understates the gap between reaching and non-reaching methods. We
therefore lead with a censoring-aware summary --- the fraction of seeds that reach
target, $\mathrm{Reach}=\tfrac15\sum_{s}\mathbf{1}[\mathrm{SR}(\theta_B)\ge0.90]$, and
the \emph{median} $Q$ among reaching seeds only. Cells still show mean$\pm$std $Q$,
but a cell is bold on $Q$ only among methods with the highest $\mathrm{Reach}$ (a
censored cell is never bold), and the headline comparison uses Reach first, then
median-$Q$ among reachers. \textbf{Eff}: demonstration efficiency, the area under the
SR-vs-\#demonstrations curve (Eq.~\ref{eq:eff}), well-defined even for non-reaching
runs. \textbf{Cov}: coverage of collected demonstrations over the task space
(Eq.~\ref{eq:cov}).

\paragraph{Seeds and significance.} We run \textbf{$5$ seeds} per cell (per-run seed
$=\text{base}+\text{run\_id}$), a floor we will raise to $10$ for the camera-ready if
resources permit. Because arms share the bootstrap per seed, seeds are \emph{paired},
so for each headline comparison we report the five paired per-seed differences
directly rather than a bootstrap CI (pinned to min/max over five points). We call a
comparison directional only when all five differences share a sign (a one-sided sign
test at $p=1/32$), paired with the effect size (median paired difference); we do
\emph{not} treat non-overlap of $\pm$std bars as significance. Given the many pairwise
comparisons, we confine strong claims to the single pre-registered contrast (PACE
vs.\ the strongest arm per cell) and report the rest descriptively. Formally, over a
fixed evaluation set $\mathcal{E}$,
\begin{equation}
\mathrm{SR}(\theta)=\frac{1}{|\mathcal{E}|}\sum_{e\in\mathcal{E}}\mathbf{1}\big[\ e\ \text{succeeds under }\pi_\theta\ \big],
\qquad \mathrm{SR}_{\mathrm{target}}=0.90,
\label{eq:sr}
\end{equation}
\begin{equation}
q(\mathrm{SR}_{\mathrm{target}})=\min\big\{\,q\ :\ \mathrm{SR}(\theta_q)\ge\mathrm{SR}_{\mathrm{target}}\,\big\},
\qquad q\leftarrow B\ \text{if never reached},
\label{eq:queries}
\end{equation}
\begin{equation}
\mathrm{Eff}=\int_{0}^{B}\mathrm{SR}(\theta_{q})\ \mathrm{d}q
\ \approx\ \sum_{q=1}^{B}\tfrac12\big(\mathrm{SR}(\theta_{q-1})+\mathrm{SR}(\theta_{q})\big),
\label{eq:eff}
\end{equation}
\begin{equation}
\mathrm{Cov}=\frac{1}{|\mathcal{Z}|}\Big|\ \big\{\,z\in\mathcal{Z}\ :\ \exists\,(s,a)\in\mathcal{D},\ \mathrm{cell}(s)=z\,\big\}\ \Big|,
\label{eq:cov}
\end{equation}
where $\mathcal{Z}$ is the quantized task space (grid cells for T1, workspace/pose
cells for T2); $\mathrm{Cov}$ isolates the effect of PACE's diversity/memory
choice on the collected data.

\subsection{Main Results}

Table~\ref{tab:main} consolidates held-out SR and query count $Q$ for all methods
across both modalities on the two implemented tasks (T1, T2). These cells answer
whether PACE reaches target with the fewest expert queries while matching or exceeding
baseline SR: averaged over the two tasks PACE attains $\PH{mean-pace-sr}$ SR, and on
Push its expert-query change to reach $0.90$ SR relative to Diff-DAgger is
$\PH{pace-vs-diff-q-reduction}$. The pre-registered contrast is PACE against \emph{the
strongest of} the DAgger family and the two single-axis controls ($R$, $W$): if PACE
does not beat both single-axis arms, the integration claim is not supported and we say
so. We report the sign from the paired per-seed differences (\S Metrics) and claim
``fewer than the strongest arm'' only when they agree in sign; with $n{=}5$ every
entry is provisional. The foundation-model and reset cost PACE alone incurs is
accounted separately in Table~\ref{tab:cost}, so ``query efficiency'' is never read in
isolation.

\begin{table}[t]
\centering
\small
\caption{Main results on the two \emph{implemented} tasks: held-out SR (\%) and
expert queries $Q$ to reach $0.90$ SR, mean$\pm$std over $5$ seeds. Best per
column within a task in \textbf{bold}, restricted to methods with the highest
Reach (a right-censored $Q$ --- target never reached, so $Q{=}B$ --- is never
bold; see \S Metrics). State/Image $=$ the two modalities (T1: equivariant MLP /
plain CNN; T2: state / image diffusion policy). Diff-DAgger is robot-only.
Reset-Curric.\ ($R$) is the single-axis \emph{where} control and IWR-Reweight ($W$)
the single-axis \emph{which} control (\S Baselines). The
planned RoboSuite tasks (T3--T5) are \emph{not} shown: they are unimplemented and
contribute no rows. All cells are red placeholders pending the five-seed runs.}
\label{tab:main}
\begin{tabular}{l cc cc}
\toprule
 & \multicolumn{2}{c}{State} & \multicolumn{2}{c}{Image} \\
\cmidrule(lr){2-3}\cmidrule(lr){4-5}
Method & SR$\uparrow$ & Q$\downarrow$ & SR$\uparrow$ & Q$\downarrow$ \\
\midrule
\multicolumn{5}{l}{\emph{T1 Toy grid}}\\
SafeDAgger     & \PH{toy-st-safe-sr}     & \PH{toy-st-safe-q}     & \PH{toy-im-safe-sr}     & \PH{toy-im-safe-q} \\
DropoutDAgger  & \PH{toy-st-dropout-sr}  & \PH{toy-st-dropout-q}  & \PH{toy-im-dropout-sr}  & \PH{toy-im-dropout-q} \\
EnsembleDAgger & \PH{toy-st-ensemble-sr} & \PH{toy-st-ensemble-q} & \PH{toy-im-ensemble-sr} & \PH{toy-im-ensemble-q} \\
ThriftyDAgger  & \PH{toy-st-thrifty-sr}  & \PH{toy-st-thrifty-q}  & \PH{toy-im-thrifty-sr}  & \PH{toy-im-thrifty-q} \\
Random         & \PH{toy-st-rand-sr}     & \PH{toy-st-rand-q}     & \PH{toy-im-rand-sr}     & \PH{toy-im-rand-q} \\
Reset-Curric.\ ($R$) & \PH{toy-st-rcur-sr} & \PH{toy-st-rcur-q} & \PH{toy-im-rcur-sr} & \PH{toy-im-rcur-q} \\
IWR-Reweight ($W$)   & \PH{toy-st-iwr-sr}  & \PH{toy-st-iwr-q}  & \PH{toy-im-iwr-sr}  & \PH{toy-im-iwr-q} \\
\textbf{PACE (ours)} & \PH{toy-st-pace-sr} & \PH{toy-st-pace-q} & \PH{toy-im-pace-sr} & \PH{toy-im-pace-q} \\
\midrule
\multicolumn{5}{l}{\emph{T2 Push (ManiSkill PushT)}}\\
Diff-DAgger    & \PH{push-st-diff-sr}     & \PH{push-st-diff-q}     & \PH{push-im-diff-sr}     & \PH{push-im-diff-q} \\
SafeDAgger     & \PH{push-st-safe-sr}     & \PH{push-st-safe-q}     & \PH{push-im-safe-sr}     & \PH{push-im-safe-q} \\
DropoutDAgger  & \PH{push-st-dropout-sr}  & \PH{push-st-dropout-q}  & \PH{push-im-dropout-sr}  & \PH{push-im-dropout-q} \\
EnsembleDAgger & \PH{push-st-ensemble-sr} & \PH{push-st-ensemble-q} & \PH{push-im-ensemble-sr} & \PH{push-im-ensemble-q} \\
ThriftyDAgger  & \PH{push-st-thrifty-sr}  & \PH{push-st-thrifty-q}  & \PH{push-im-thrifty-sr}  & \PH{push-im-thrifty-q} \\
Random         & \PH{push-st-rand-sr}     & \PH{push-st-rand-q}     & \PH{push-im-rand-sr}     & \PH{push-im-rand-q} \\
Reset-Curric.\ ($R$) & \PH{push-st-rcur-sr} & \PH{push-st-rcur-q} & \PH{push-im-rcur-sr} & \PH{push-im-rcur-q} \\
IWR-Reweight ($W$)   & \PH{push-st-iwr-sr}  & \PH{push-st-iwr-q}  & \PH{push-im-iwr-sr}  & \PH{push-im-iwr-q} \\
\textbf{PACE (ours)} & \PH{push-st-pace-sr} & \PH{push-st-pace-q} & \PH{push-im-pace-sr} & \PH{push-im-pace-q} \\
\bottomrule
\end{tabular}
\end{table}

\subsection{Cost Accounting: What $Q$ Does and Does Not Count}

The query metric $Q$ counts \emph{only} expert demonstrations, which the
one-demo-per-round budget equalizes across methods. But PACE's Execute stage is not
free along two axes the baselines never touch. \emph{(1) Reset access} (only under the
optional prescription): the on-policy default resumes the expert on the failed
rollout, but a \texttt{BRIDGE} prescription instead samples a new start state
$s_0\sim\Xi(\cdot\mid\xi)$ (Eq.~\ref{eq:prescribe-demo}), requiring a resettable
simulator (Limitations~ii). \emph{(2) Foundation-model calls}: each PACE round issues
VLM/LLM calls the baselines do not. Table~\ref{tab:cost} reports these axes per
reached-target run so PACE's $Q$ advantage is never read without its reset and
inference cost. Because a single $-$Execute row would confound losing the prescribed
start state with losing reset access, we report both controls in
Table~\ref{tab:ablation} ($-$Execute and $-$Prescribe$+$Reset), letting a reviewer
attribute any drop to the right cause.

\begin{table}[t]
\centering
\small
\caption{Competing cost axes that the query count $Q$ omits, per run that reaches
$0.90$ SR (T2 Push, mean over $5$ seeds). Baselines require none of the last three
columns. All cells are red placeholders pending the runs.}
\label{tab:cost}
\begin{tabular}{l cccc}
\toprule
Method & $Q\downarrow$ & Resets & VLM calls & LLM calls \\
\midrule
Diff-DAgger          & \PH{push-diff-q}      & $0$ & $0$ & $0$ \\
Best baseline        & \PH{push-base-q}      & $0$ & $0$ & $0$ \\
\textbf{PACE (ours)} & \PH{push-pace-q}      & \PH{pace-resets} & \PH{pace-vlm-calls} & \PH{pace-llm-calls} \\
\bottomrule
\end{tabular}
\end{table}

\subsection{Learning Curves and Coverage}

Endpoint SR and Q answer \emph{whether} PACE is more query-efficient; the curves and
coverage test \emph{why}. We plot held-out SR against the number of collected
demonstrations for every implemented (task, modality) cell (Figure~\ref{fig:lc}). The
question is whether PACE lies on the upper-left of the SR-vs-demonstration frontier
--- higher \textbf{Eff} at a given budget --- because cross-round memory rotation
front-loads failure-mode coverage rather than re-correcting the dominant mode; whether
it holds is decided by the pending five-seed numbers. On the state PushT curve, PACE
attains $\PH{push-st-pace-eff}$ Eff vs.\ $\PH{push-st-diff-eff}$ for Diff-DAgger; on
the image curve, $\PH{push-im-pace-eff}$ vs.\ $\PH{push-im-diff-eff}$. Coverage
\textbf{Cov} is the mechanistic check: PACE reaches $\PH{push-st-pace-cov}$
workspace-cell coverage on state PushT (vs.\ $\PH{push-st-diff-cov}$ for the strongest
baseline) and $\PH{toy-st-pace-cov}$ grid-cell coverage on the toy task (vs.\
$\PH{toy-st-rand-cov}$); if the efficiency ordering tracks the coverage ordering, the
gain is attributable to broader failure-mode coverage rather than higher
per-demonstration quality. Figure~\ref{fig:qual} shows a representative
failure$\to$prescription panel: a clustered failure mode, its VLM-perceived
description, and the prescribed corrective reset state.

\begin{figure}[t]
\centering
\includegraphics[width=0.95\columnwidth]{figures/lc_push_image.pdf}
\caption{Learning curves: held-out success rate versus number of collected
demonstrations (five-seed mean; shaded band is $\pm1$ std), shown for a
representative (task, modality) cell. PACE is expected to lie on the upper-left
frontier --- reaching $\mathrm{SR}_{\mathrm{target}}=0.90$ (dashed) with fewer
demonstrations --- if the cross-round memory rotation front-loads
failure-mode coverage. One panel per task$\times$modality appears in the full grid;
per-panel demo-efficiency (Eq.~\ref{eq:eff}) is reported in the main tables.}
\label{fig:lc}
\end{figure}

\begin{figure}[t]
\centering
\includegraphics[width=0.95\columnwidth]{figures/qual_push.pdf}
\caption{Qualitative failure$\to$prescription example. Left: a clustered failure
mode (members of $C_{\mathrm{tgt}}$ at their peak diffusion-loss frames,
Eq.~\ref{eq:perceive}). Center: the VLM-perceived description grounding the scene
command. Right: the prescribed, workspace-clamped corrective reset state $\xi$
(Eq.~\ref{eq:prescribe}) from which the single new expert demonstration is
collected.}
\label{fig:qual}
\end{figure}

\subsection{Ablations}

To attribute the gains to the PACE mechanisms (Q4) we ablate one at a time
(Table~\ref{tab:ablation}). \textbf{The analysis is pre-registered; it is not an
AND-rule.} We run the ablation on \emph{two} cells --- T1 Toy (state) and T2 Push
(image), spanning both embodiments and modalities --- and report each cell separately
with its paired-per-seed effect size (median paired difference from full PACE), not a
strict ``degrades in both'' gate that would discard true effects as readily as false
ones over two noisy $n{=}5$ cells. At $n{=}5$ the paired sign test resolves a
consistent-direction effect at $p{=}1/32$, and the median paired difference
distinguishes a genuine null from low power. A mechanism ``carries weight'' in a cell
when its removal degrades that cell with all five differences agreeing in sign; we
state the per-cell direction expected \emph{a priori}. Each row isolates one
mechanism (defined precisely in the table caption), separating loss-localization
($-$Perceive) from the foundation models ($-$VLM, $-$LLM, which together answer
whether these are causal or decorative), the failure-mode partition ($-$Assess), the
coverage mechanism ($-$Memory, Eq.~\ref{eq:memory}), and the within-round context
($-$Choose-ctx). On the Execute side, the ($-$Execute,$-$Prescribe$+$Reset,full)
triple separates free resets from the prescribed start state. We expect the largest
Eff drop at $-$Memory or $-$Assess and, per our stated scope, a small or null
$-$Prescribe$+$Reset$\to$full gap; if that gap is null, the prescription option is not
a driver and we say so. We also report the $\kappa$ sweep and the Qwen3
LLM~\cite{yang2025qwen3} reasoning-enabled vs.\ disabled modes ($\PH{abl-think-eff}$
vs.\ $\PH{abl-nothink-eff}$), with failure perception driven by the Qwen3-VL
vision--language model~\cite{bai2025qwen3vl}.

\begin{table}[t]
\centering
\small
\caption{Ablation of PACE mechanisms on two cells (T1 Toy state, T2 Push image),
mean$\pm$std over $5$ seeds, reported \emph{per cell} (no both-cells AND-rule; see
text). $-$Perceive ablates loss-localization only; $-$VLM/$-$LLM remove the
foundation models; $-$Memory removes the acting coverage mechanism; $-$Choose-ctx
($\kappa{=}1$) removes the within-round \texttt{BRIDGE} context and is expected
\emph{near-identical} to full PACE under the on-policy default; $-$Execute is fully
on-policy with no reset access; $-$Prescribe$+$Reset keeps reset access but does not
move the start state, so ($-$Execute,$-$Prescribe$+$Reset,full) separates free
resets from the prescribed \emph{where}. Demonstration
efficiency Eff (Eq.~\ref{eq:eff}), coverage Cov (Eq.~\ref{eq:cov}). All cells are
red placeholders. On T1 the VLM/LLM operate on the grid corridor prescription;
$-$VLM/$-$LLM there use the scripted A$^\star$ corridor.}
\label{tab:ablation}
\begin{tabular}{l cc cc}
\toprule
 & \multicolumn{2}{c}{T1 Toy (state)} & \multicolumn{2}{c}{T2 Push (image)} \\
\cmidrule(lr){2-3}\cmidrule(lr){4-5}
Variant & Eff$\uparrow$ & Cov$\uparrow$ & Eff$\uparrow$ & Cov$\uparrow$ \\
\midrule
$-$Perceive & \PH{abl-toy-perceive-eff} & \PH{abl-toy-perceive-cov} & \PH{abl-perceive-eff} & \PH{abl-perceive-cov} \\
$-$VLM      & \PH{abl-toy-vlm-eff}      & \PH{abl-toy-vlm-cov}      & \PH{abl-vlm-eff}      & \PH{abl-vlm-cov} \\
$-$LLM      & \PH{abl-toy-llm-eff}      & \PH{abl-toy-llm-cov}      & \PH{abl-llm-eff}      & \PH{abl-llm-cov} \\
$-$Assess   & \PH{abl-toy-assess-eff}   & \PH{abl-toy-assess-cov}   & \PH{abl-assess-eff}   & \PH{abl-assess-cov} \\
$-$Memory   & \PH{abl-toy-memory-eff}   & \PH{abl-toy-memory-cov}   & \PH{abl-memory-eff}   & \PH{abl-memory-cov} \\
$-$Choose-ctx ($\kappa{=}1$) & \PH{abl-toy-choose-eff} & \PH{abl-toy-choose-cov} & \PH{abl-choose-eff} & \PH{abl-choose-cov} \\
$-$Execute (no reset) & \PH{abl-toy-execute-eff}  & \PH{abl-toy-execute-cov}  & \PH{abl-execute-eff}  & \PH{abl-execute-cov} \\
$-$Prescribe$+$Reset  & \PH{abl-toy-presreset-eff} & \PH{abl-toy-presreset-cov} & \PH{abl-presreset-eff} & \PH{abl-presreset-cov} \\
\midrule
\textbf{PACE (full)} & \PH{abl-toy-full-eff} & \PH{abl-toy-full-cov} & \PH{abl-full-eff} & \PH{abl-full-cov} \\
\bottomrule
\end{tabular}
\end{table}

\section{Conclusion}
\label{sec:conclusion}

\textbf{PACE} changes what the DAgger family asks. Safe and query-efficient DAgger
variants decide only \emph{when} to query the expert from a per-state uncertainty
or safety
gate \cite{zhang2017safedagger,menda2017dropoutdagger,menda2019ensembledagger,hoque2021thriftydagger,lee2025diffdagger};
PACE adds \emph{which} failure mode to correct --- its central axis --- and an
\emph{optional} decision of \emph{where} the corrective demonstration begins. Failures
are localized at their peak-loss frame, optionally described by a vision--language model
\cite{bai2025qwen3vl}, and clustered into
failure modes (Eqs.~\eqref{eq:perceive}--\eqref{eq:partition}); the target mode
is chosen by a cross-round memory rotation with a small k-center context set
(Eqs.~\eqref{eq:memory}--\eqref{eq:diversity}); the demonstration is collected
on-policy by default, or, when enabled, from an LLM-prescribed
\cite{yang2025qwen3} sub-task-entry start state
(Eqs.~\eqref{eq:prescribe}--\eqref{eq:prescribe-demo}). As a unifying view scoped to
the diffusion-loss task rather than a theorem, prior interactive IL sits at one
corner: collapsing selection, memory, and the prescribe map turns a PACE round
into a single on-policy query, and at window $K{=}1$ that corner meets the
diffusion-loss rule of Diff-DAgger \cite{lee2025diffdagger} under the
peak-loss/first-crossing monotonicity condition noted in
\S\ref{sec:method}. We claim nothing stronger, and the view does not cover the toy
grid's classifier loss.

On the two implemented tasks of Table~\ref{tab:tasks}---a toy $5\times5$ grid and ManiSkill
Push \cite{florence2021implicitbc,tao2024maniskill3}, in both state and image modalities---we
test whether PACE reaches the $0.90$ target success rate with
fewer expert queries than the DAgger-family baselines while matching or exceeding their
final success rate, reporting paired per-seed differences and
accounting separately for reset and foundation-model cost (\S Metrics,
Table~\ref{tab:cost}). \emph{The result cells (mean held-out SR \PH{mean-pace-sr}, query
change \PH{pace-vs-diff-q-reduction} relative to Diff-DAgger, Table~\ref{tab:main}) are
placeholders; no numeric claim is yet established until the five-seed runs complete.}
Three RoboSuite tasks on a UR5/UR5e arm \cite{zhu2020robosuite}
are a \emph{planned} cross-embodiment extension, not evaluated here. The
ablation (Table~\ref{tab:ablation}, with VLM-off, LLM-off, and reset-access controls)
tests whether the \emph{which} mechanisms carry the gain and whether the optional
\emph{where} adds anything on top. Because integration-only novelty is weak, the
headline contrast pits PACE against the strongest DAgger baseline and the two
single-axis controls. We expect the signature to appear in coverage (Eq.~\eqref{eq:cov}),
consistent with the guiding idea that spending expert effort on curated, non-redundant
failure modes---rather than on every uncertain state---yields more sample-efficient
corrective data. What we defend today is the method, framework, and protocol;
the query-efficiency thesis awaits the five-seed T1/T2 numbers.

\subsection{Limitations}
\label{sec:limitations}

We make the load-bearing assumptions explicit. \emph{(i) Evaluation scope.} The
study covers two implemented tasks (T1 Toy, T2 Push) at five seeds; the three
RoboSuite/UR5 tasks are specified but not run, all result cells are placeholders,
and $n{=}5$ limits statistical resolution (\S Metrics).
\emph{(ii) Reset access and cost.} The optional \emph{where} prescription is
\emph{conditional} on a resettable simulator
\cite{eysenbach2018leavenotrace,florensa2017reversecurriculum}, and every round
issues VLM/LLM calls; the query metric $q$ counts neither, so we report them as a
separate cost axis (Table~\ref{tab:cost}) and two controls that
separate free resets from the prescribed start state (Table~\ref{tab:ablation})
rather than folding them into $q$. The on-policy default needs no reset access.
\emph{(iii) Foundation-model fidelity.} A mis-described failure or
out-of-bounds prescription degrades the demo; we mitigate with hard KAG workspace
clamps (Eq.~\eqref{eq:prescribe}), a non-empty floor, and a deterministic
escalation ladder, but the pipeline is only as reliable as its backbones and
prompts, and whether the VLM is causal is left to the $-$VLM ablation.
\emph{(iv) Privileged experts and hand-designed features.} Our expert is an
A$^\star$/BFS or PPO / motion-planner oracle; a noisier human expert may differ
\cite{mandlekar2020iwr,liu2023sirius}, and Assess uses a fixed $6$-D geometric or
frozen-R3M \cite{nair2022r3m} descriptor with $\kappa$ and the memory bandwidth
fixed, not learned. \emph{(v) Baseline fidelity.} SafeDAgger$^\star$ uses an oracle
action-discrepancy gate rather than a learned safety classifier, a variant whose
bias direction is \emph{unknown} (\S Baselines),
and on the robot path ThriftyDAgger's risk arm degrades to novelty-only when its
success-Q buffer cannot be read; both are faithful-as-possible proxies whose
deviation we do not claim to bound.

\subsection{Future Work}
\label{sec:future-work}

The immediate step is completing the T1--T2 five-seed runs (and raising to ten
seeds) and the planned RoboSuite tasks, so the design can be judged on results.
Several extensions then open. One is to close the loop between Perceive and
Assess: letting the VLM's semantic failure taxonomy \emph{define} the clustering
feature space --- rather than clustering hand-designed or frozen-encoder features
--- could make mode discovery more robust and give the VLM a say in \emph{which}
failures are grouped
\cite{liu2023reflect,duan2025aha,ash2020badge,sener2018coreset,settles2009active}.
A second would relax resettability, growing Execute from a single
reset pose into on-hardware sub-task curricula run with human experts
\cite{kelly2019hgdagger,liu2023sirius}. And because
PACE is agnostic to the learner, it should carry over to other generative policy
classes --- trajectory diffusion planners \cite{janner2022diffuser} and
vision--language--action models \cite{brohan2023rt2,driess2023palme} --- pushing
failure-driven data collection past the diffusion policy studied here
\cite{chi2023diffusionpolicy}. What ties these together is a change of
framing: treat interactive imitation learning as \emph{failure-mode curation}. An
agent that groups its own mistakes and corrects them can put a scarce expert
budget where it is needed most.

\appendix
\section{Appendix: Baseline Query Predicates}
\label{app:baselines}

For completeness we give the five query predicates deferred from
\S\ref{sec:unified}. Each is the score inside the unified template
$\mathrm{Query}_{\mathrm{method}}(s_t)=\mathbf{1}[\mathrm{score}_{\mathrm{method}}(s_t)\gtrless\tau_{\mathrm{method}}]$,
with all action comparisons taken in the shared relative-joint-position space.
\emph{SafeDAgger}$^\star$ \cite{zhang2017safedagger} gates on novice--expert action
discrepancy (on the grid, the score is the fraction of off-optimal-mask steps),
\begin{equation}
\mathrm{Query}_{\mathrm{Safe}}(s_t)=\mathbf{1}\big[\ \|a^{\mathrm{nov}}_t-a^{\mathrm{exp}}_t\|_2 > \tau_{\mathrm{sd}}\ \big].
\label{eq:safe}
\end{equation}
\emph{DropoutDAgger} \cite{menda2017dropoutdagger,gal2016dropout} draws
$N_{\mathrm{mc}}$ stochastic samples and queries when the in-$\tau$-ball agreement
$\hat p_t$ with the expert drops below $p$,
\begin{align}
\hat p_t=\tfrac{1}{N_{\mathrm{mc}}}\sum_{i=1}^{N_{\mathrm{mc}}}\mathbf{1}\big[\ \|a^{(i)}_t-a^{\mathrm{exp}}_t\|_2\le\tau\ \big],
\qquad
\mathrm{Query}_{\mathrm{Drop}}(s_t)=\mathbf{1}\big[\ \hat p_t<p\ \big].
\label{eq:dropout}
\end{align}
\emph{EnsembleDAgger} \cite{menda2019ensembledagger,lakshminarayanan2017ensembles}
queries on ensemble doubt (variance) \emph{or} mean discrepancy over $M$ members,
\begin{align}
\mathrm{doubt}_t&=\tfrac1{d_a}\!\sum_{j=1}^{d_a}\mathrm{Var}_{m}\big(a^{[m]}_{t,j}\big),
\quad
\mathrm{disc}_t=\tfrac1M\!\sum_{m=1}^{M}\big\|a^{[m]}_t-a^{\mathrm{exp}}_t\big\|_2, \nonumber\\
\mathrm{Query}_{\mathrm{Ens}}(s_t)&=\mathbf{1}\big[\ \mathrm{disc}_t>\tau\ \ \lor\ \ \mathrm{doubt}_t>\chi\ \big].
\label{eq:ensemble}
\end{align}
\emph{ThriftyDAgger} \cite{hoque2021thriftydagger} queries on novelty (doubt)
\emph{or} task risk $1-Q$. The success-Q value $Q_\psi$ is trained by BCE on a
Monte-Carlo success-to-go target $y_t$, where $n$ is the episode length (the number
of steps in the trajectory the pair is drawn from); the thresholds
$(\delta_h,\beta_h)$ are the $(1-\alpha_h)$-quantiles of the current round's
per-state doubt/risk pool (the calibration set),
\begin{align}
y_t&=\mathbf{1}[\mathrm{success}]\cdot\gamma^{\,n-1-t},
\quad
\mathcal{L}_Q=\mathrm{BCE}\big(\sigma(Q_\psi(s_t,a_t)),\,y_t\big),
\quad
\mathrm{risk}(s,a)=1-\sigma\!\big(Q_\psi(s,a)\big), \nonumber\\
\delta_h&=\mathrm{Quantile}_{1-\alpha_h}\{\mathrm{doubt}\},\
\beta_h=\mathrm{Quantile}_{1-\alpha_h}\{\mathrm{risk}\},\
\mathrm{Query}_{\mathrm{Thr}}(s_t)=\mathbf{1}\big[\ \mathrm{doubt}_t>\delta_h\ \lor\ \mathrm{risk}_t>\beta_h\ \big].
\label{eq:thrifty}
\end{align}
The uniform-random control (\emph{Random}) queries at a pre-drawn step, a lower
bound on any informed rule,
\begin{equation}
t_{\mathrm{stag}}\sim\mathrm{Unif}\{0,\dots,H-1\},
\qquad
\mathrm{Query}_{\mathrm{Rand}}(s_t)=\mathbf{1}\big[\ t\ge t_{\mathrm{stag}}\ \big].
\label{eq:stagger}
\end{equation}
Threshold constants $(\tau_{\mathrm{sd}},\tau,p,\chi,\alpha_h,\alpha,K,M,N_{\mathrm{mc}})$
are fixed configuration values (defaults in \S Baselines) rather than tuned per task.
As noted in \S\ref{sec:unified}, SafeDAgger$^\star$ substitutes an oracle
action-discrepancy signal for the original learned safety classifier (a conservative
proxy), and on the robot path ThriftyDAgger degrades to novelty-only when its
success-Q buffer is unavailable.

\section*{Appendix: T1 Grid-World Instantiation of Perceive}
\label{app:t1grid}

Equations~\ref{eq:perceive}--\ref{eq:prescribe} are written for the manipulation
setting; the toy grid (T1) uses a distinct but explicitly-defined instantiation of
the same pipeline, and the clustering, dominance, memory, and diversity stages
(Eqs.~\ref{eq:partition}--\ref{eq:diversity}) run unchanged on the resulting $6$-D
grid descriptor. There is no continuous yaw, TCP, or object pose, so we set
$p^{\mathrm{tee}}_i=(\mathrm{col}_i,\mathrm{row}_i)$ to the agent's grid cell at the
failure step, replace the sin/cos-yaw pair with the signed cell offset to the goal
$(g_x-\mathrm{col}_i,\,g_y-\mathrm{row}_i)$, keep $\rho_i=t^\star_i/\max(1,T_i)$, and
set $\delta_i=\|(\mathrm{col}_i,\mathrm{row}_i)-(g_x,g_y)\|_1$ (Manhattan
distance-to-goal, the grid analogue of contact distance). The per-step signal
$\ell_t$ is the classifier's behavior-cloning loss rather than the diffusion loss,
and the peak-loss step $t^\star_i=\arg\max_t\ell_t$ localizes the failure exactly as
in Eq.~\ref{eq:perceive}. The prescribe map $g$ (Eq.~\ref{eq:prescribe}) is
likewise replaced on the grid by a \emph{corridor} prescription: the LLM emits a
start/goal/fire layout and a cell-sequence corridor, which an A$^\star$ expert is
forced down to yield the one new demonstration. The anchor-relative perturbation
caps and workspace clamp of Eq.~\ref{eq:prescribe} become an A$^\star$/BFS
feasibility check that rejects corridors leaving the grid, stepping on fire, or
disconnecting start from goal. As noted in the Method, the corner correspondence to
the query family is vacuous here (the per-step signal is a classifier BCE loss
rather than a diffusion loss); it holds only at the informal level of
``highest-loss failure, corrected on-policy.''

\bibliographystyle{aaai2027}
\bibliography{references}

\end{document}
